% Philippica XIII (philippic-13) --- English translation
% Translated by Claude (Anthropic) from the Latin (Perseus).
% Style guide: STYLE.md.

\ciceroSection{1} From the very beginning of this war, senators --- the war we took up against impious and criminal citizens --- I feared that some treacherous offer of peace might quench our zeal to recover our liberty. For even the name of peace is sweet, and the thing itself is at once pleasant and wholesome. Yet a man who delights in discord, in the slaughter of citizens, in civil war seems to hold dear neither his own hearth, nor the public laws, nor the rights of liberty; and such a man I judge fit to be cast out of the number of human beings, banished from the bounds of human nature. And so, whether it was Sulla or Marius or both of them, whether Octavius or Cinna or Sulla a second time, or the second Marius and Carbo, or any other who longed for civil war, I judge him a citizen to be detested, born to be the bane of the commonwealth.

\ciceroSection{2} For what shall I say of the most recent case --- the man whose acts we defend, while we confess that the author of them was justly slain? Nothing, then, is more loathsome than this citizen, nothing than this man --- if he is to be counted a citizen or a man at all --- who craves civil war. But this, senators, must first be examined: whether peace is possible with all men, or whether there is some war past atonement, in which the bargain of peace would be the charter of slavery. When Sulla made peace with Scipio --- or pretended to --- there was no cause to despair that, if it held, the state might reach some tolerable condition. Had Cinna been willing to confirm a concord with Octavius, sanity might have remained among men in public life. In the most recent war, had Pompey relented somewhat from his utter sternness, and Caesar much from his ambition, we might have kept both a stable peace and some commonwealth still. But this case --- what is it? Can there be peace with the Antonii? With Censorinus, Ventidius, Trebellius, Bestia, Nucula, Munatius, Lento, Saxa? I have named only a few by way of example: the breed is past numbering, and the savagery of the rest you can see for yourselves.

\ciceroSection{3} Add to these the wrecks among Caesar's friends --- your Barbae and Cassii, the Barbatii, the Pollios; add Antony's playmates and cronies, Eutrapelus, Mela, Pontius, Caelius, Crassicius, Tiro, Mustela, Petusius. The retinue I leave aside; I name only the captains. To these are joined the Larks and the rest of the veterans, that nursery of jurors for the third decury, who, having drained their own fortunes and gorged on Caesar's bounty, have set their hearts on ours.

\ciceroSection{4} O the faithful right hand of Antony, with which he has butchered so many of his fellow citizens! O the firm and holy treaty we shall have made with the Antonii! If Marcus tries to violate it, the sanctity of Lucius will call him back from the crime. If there is room for them in this city, there will be no room for the city itself. Set their faces before your eyes --- above all the faces of the Antonii: the gait, the look, the countenance, the swagger; some friends guarding their flanks, others marching before them. What reek of wine, what insults, what threats of words do you suppose there will be! --- unless, perhaps, peace itself will soften them, and, once they have come into this house, they will greet us graciously and address each of us by name with courtesy.

\ciceroSection{5} Do you not remember, by the immortal gods, the votes you cast against them? You annulled the acts of Marcus Antonius; you took down his laws; you declared that they had been carried by violence and against the auspices; you stirred up levies throughout all Italy; you judged his colleague --- the partner of his every crime --- a public enemy. What peace can there be with such a man? If the enemy were a foreign one, even that would scarcely be possible after deeds like these; but somehow it might be. Seas, mountains, great stretches of country would lie between; you would hate one you did not have to see. But these men will cling before our eyes, and, when they may, at our throats. With what fences shall we pen in beasts so monstrous? But the issue of war, you say, is uncertain. It belongs wholly to brave men --- which is what you ought to be --- to make good their valour (for that much lies in their power) and not to dread the blame that belongs to fortune.

\ciceroSection{6} But since this order is asked to furnish not courage only but wisdom as well --- though the two can scarcely be parted, let us part them all the same --- courage bids us fight, kindles a just hatred, drives us to the clash, summons us to danger: and wisdom? She follows more cautious counsels, looks to the future, is on every side more guarded. What, then, does she advise? For we must obey her, and judge that best which has been most wisely determined. If she charges me to count nothing dearer than life, to hazard no peril to my head, to flee every crisis, I will ask of her: `And even so --- when I have done all this --- must I then be a slave?' If she nods yes, then for my part, learned though this wisdom be, I will not listen to it. But if she answers, `Guard your life and body, your fortunes and your estate, yet in such a way that you reckon them all of less account than liberty; resolve to enjoy them only if you may with the commonwealth free; and do not cast away liberty for their sake, but cast them away for liberty's sake, as pledges forfeit to wrong' --- then I shall think I am hearing the voice of wisdom, and I shall obey her as a god.

\ciceroSection{7} And so, if with these men received back among us we can still be free, let us conquer our hatred and put up with peace; but if there can be no quiet for us while they are unharmed, let us rejoice that the chance to fight it out has been offered. For either, when they are killed, we shall enjoy a victorious commonwealth, or, if we are crushed --- may Jupiter avert the omen! --- we shall live, if not in our breath, yet in the renown of our valour. But Marcus Lepidus, it is urged --- twice hailed imperator, pontifex maximus, who in the last civil war served the commonwealth most nobly --- exhorts us to peace. No man's authority weighs more with me, senators, than that of Marcus Lepidus, whether for his own worth or for the dignity of his house. There are added to this many great private services of his to me, and certain good offices of mine to him. But the greatest benefit of his I reckon to be this: that he is of such a mind toward the commonwealth, which has always been dearer to me than my own life.

\ciceroSection{8} For when by his authority he led Magnus Pompeius --- a young man of the highest distinction, the son of a most excellent father --- to peace, and freed the commonwealth without arms from a great peril of civil war, then I judged myself bound by his service beyond what one man's share requires. And so I both decreed him the most ample honours I could, in which you concurred with me, and never ceased to hope and to speak the best of him. By many great pledges the commonwealth holds Marcus Lepidus bound to her. His nobility is of the highest; he has every office, the most exalted priesthood, and many ornaments of the city --- monuments of his own, his brother's, his ancestors'; a most admirable wife, most cherished children, an estate at once ample and unstained by the blood of citizens. No citizen has been harmed by him; many have been saved by his kindness and his mercy. Such a man, then, and such a citizen, may err in his opinion, but he cannot by any means dissent in will from the commonwealth. Marcus Lepidus wishes for peace.

\ciceroSection{9} Splendidly so --- if he can bring about such a peace as he lately brought about: a peace by which the commonwealth shall look once more upon the son of Gnaeus Pompeius and take him back into her bosom and embrace, and shall think that not he alone, but she herself with him, has been restored to her. This was the reason you decreed him a statue on the Rostra with a glorious inscription, and a triumph in his absence. For though he had waged great wars, and wars deserving of a triumph, that was not what was granted him --- the like was given neither to Lucius Aemilius, nor to Aemilianus Scipio, nor to the elder Africanus, nor to Marius, nor to Pompey, who waged greater wars; rather, because he had brought a civil war to its close in silence, you conferred on him, the moment it was permitted, the greatest honours of all.

\ciceroSection{10} Do you suppose, then, Marcus Lepidus, that the Antonii will be such citizens in the commonwealth as the Pompeius she is to have? In the one, modesty, gravity, restraint, integrity; in the others --- and when I challenge these, I pass over in my mind not a single man of that bandit herd --- lusts, crimes, a monstrous daring for every outrage. And then I beg you, senators: which of you does not see what Fortune herself, who is called blind, has seen? For with Caesar's acts kept intact --- the acts we defend for concord's sake --- Pompeius's own house will lie open to him, and he will buy it back for no less than Antony paid for it; he will buy it back, I say --- the son will buy back the house of Gnaeus Pompeius. A bitter thing! But this has been lamented long enough and bitterly enough. You decreed for Pompeius a sum as great as a victorious enemy had realized from his father's goods in the scattering of the spoils.

\ciceroSection{11} But this stewardship I claim for myself, in right of my old friendship and tie with his father: he will buy back the gardens, the town house, certain city properties that Antony holds. For the silver, the clothing, the furniture, the wine --- which that glutton has squandered --- he will let go with an even mind. The Alban and Formian estates he will recover from Dolabella; the Tusculan too, even from Antony; and from the Falernian let those Geese be driven off who now assault Mutina and besiege Decimus Brutus. There are perhaps more such, but they slip from my memory. I say further that even those who are not counted among the enemy shall give back to the son the Pompeian possessions at the price for which they bought them.

\ciceroSection{12} It was the act of a man rash enough --- not to say reckless --- to lay a hand on any of those goods at all; but who will be able to keep them, now that their illustrious owner is restored? Will he not give them back --- he who, coiled about his master's patrimony like a serpent about its treasure, the slave of Pompey, the freedman of Caesar, has seized the holdings of the Lucanian land? And that seven hundred million sesterces which you pledged to the young man, senators, shall be so apportioned that the son of Gnaeus Pompeius will seem to have been settled by you in his own inheritance. So much from the Senate: the rest the Roman people will pursue on behalf of that family, which it has seen in its full greatness --- and first of all his father's place in the augural college, into which I, that I may render to the son what I received from the father, will co-opt him by my own nomination. Which augur, then, will Jupiter Best and Greatest --- whose interpreters and intermediaries we are appointed to be --- which will the Roman people more gladly ratify: Pompeius or Antony? To me, indeed, it seems that by the will of the immortal gods Fortune has wished this: that, with Caesar's acts made firm and ratified, the son of Gnaeus Pompeius might recover both the dignity and the fortunes of his father.

\ciceroSection{13} Nor do I think this should be passed over in silence, senators: that the envoys --- those most distinguished men Lucius Paulus, Quintus Thermus, Gaius Fannius, whose goodwill toward the commonwealth you have found to be lasting and unwavering --- report that they turned aside to Massilia to meet Pompeius, and found him most ready in spirit to march with his forces to Mutina, did he not fear to give offence to the feelings of the veterans. But he is truly the son of that father who did no fewer things wisely than bravely. And so you understand both that the spirit was ready in him and that the judgment was not wanting. And Marcus Lepidus too must take care in this --- that he seem to do nothing more presumptuous than his own character warrants.

\ciceroSection{14} For if he means to frighten us with his army, he forgets that that army belongs to the Senate and the Roman people and the whole commonwealth, not to himself. `But he can use it as his own.' What of that? Is everything that good men have the power to do therefore to be done --- even if it will be base, even if ruinous, even if it will not be lawful at all? And what is baser, what fouler, what less becoming than to lead an army against the Senate, against the citizens, against the fatherland? And what more to be condemned than to do a thing that is not lawful? But it is lawful for no man to lead an army against the fatherland --- if indeed we call lawful only what is granted by the laws, by the usage and the institutions of our ancestors. For what each man is able to do is not therefore lawful for him, nor, if he is not stopped, is the thing thereby permitted. To you, Lepidus, as to your ancestors, the fatherland gave an army for her own sake. With it you will keep off the enemy, you will extend the bounds of empire: you will obey the Senate and the Roman people, should it perchance turn you to some other task.

\ciceroSection{15} If these are your thoughts, you are indeed Marcus Lepidus, pontifex maximus, great-grandson of Marcus Lepidus, pontifex maximus; but if you judge that men may do as much as they have the power to do, take care that you do not seem to prefer examples that are other men's, and recent, to those that are ancient and your own house's. And if you interpose your authority without arms, I praise you the more --- yet take care lest even this be unnecessary. For though your authority is as great as it ought to be in a man of the highest birth, still the Senate does not despise itself; nor was it ever more weighty, more steadfast, more brave. Fired all of us, we are swept on to recover our liberty; the ardour of so great a Senate and Roman people cannot be quenched by any man's authority; we hate, we fight in our anger, the arms cannot be wrenched from our hands; we cannot bear to hear the signal for retreat or the recall from war; we hope for the best, but we would rather endure even the hardest things than be slaves.

\ciceroSection{16} Caesar has made ready an unconquered army; two most valiant consuls are at hand with their forces; the various and great reinforcements of Lucius Plancus, the consul-designate, are not wanting; the struggle is for the safety of Decimus Brutus; one frenzied gladiator, with a band of the foulest brigands, wages war against the fatherland, against the household gods, against the altars and the hearths, against four consuls. To this man shall we yield, to his terms shall we listen, with this man shall we believe peace can be made? `But there is a danger that we may be overwhelmed.' I have no fear that one who cannot enjoy his own most ample fortune unless we are safe will betray his own safety. Good citizens nature makes first; then fortune lends her aid. For it is the interest of all good men that the commonwealth be safe; but in those who are fortunate this shows the more plainly.

\ciceroSection{17} Who is more fortunate than Lepidus, as I said before --- who, at the same time, more sane? The Roman people saw his dejection and his tears at the Lupercalia; it saw how cast down, how broken he was, when Antony, in setting the diadem upon Caesar, would rather be that man's slave than his colleague. And if he had been able to keep clear of his other shameful and criminal deeds, for this one act alone I would still think him deserving of every punishment. For if he himself could be a slave, why was he forcing a master upon us? And if his own boyhood had borne the lusts of those who played the tyrant over him, was he therefore to procure a master and a tyrant for our children too? And so, when that man was killed, just as he had wished him to be toward us, so did he himself prove toward all the rest.

\ciceroSection{18} For in what land of barbarians has there ever been a tyrant so foul, so cruel, as Antony, hedged about in this city with the arms of barbarians? Under Caesar's mastery we used to come into the Senate, if not freely, yet at least in safety. Under this arch-pirate --- for why should I call him a tyrant? --- these benches were occupied by Ituraeans. He burst out suddenly to Brundisium, to advance from there in close column upon the city; Suessa, a most splendid town --- now of most honourable burgesses, once of colonists --- he filled with the blood of the bravest soldiers; at Brundisium, in the bosom of a wife not only most greedy but most cruel, he butchered the chosen centurions of the Martian legion. From there, with what frenzy, with what fury he hurled himself upon the city --- that is, upon the slaughter of every best man! At which time the immortal gods themselves offered us a protection unforeseen and beyond our hope.

\ciceroSection{19} For the incredible and godlike valour of Caesar checked the cruel and frenzied onsets of the brigand --- whom, mad as he was, he thought to wound by his edicts, not knowing that whatever falsehoods he uttered against that most blameless young man recoiled in truth upon the memory of his own boyhood. With what an escort, or rather what a column, he entered the city! --- threatening masters to right and left, while the Roman people groaned, marking down houses, openly promising his men that he would parcel out the city among them. And on that very day he passed countless decrees of the Senate, all of which, indeed, were delivered sooner than they were written. He returned to the soldiers; then came that pestilent harangue at Tibur. From there a rush to the city; the Senate to the Capitol; a consular motion was ready for hemming in the young man, when suddenly --- for he knew that the Martian legion had encamped at Alba --- news of the Fourth was brought to him. Stunned by it, he threw aside his plan of putting Caesar's case before the Senate; he went out, not by the highways but by the byways, in his general's cloak.

\ciceroSection{20} From that moment it was no march, but a running flight into Gaul. He supposed that Caesar was at his heels with the Martian legion, with the Fourth, with the veterans, whose very name for fear he could not endure; and as he was pushing into Gaul, Decimus Brutus threw himself in his path --- choosing to be encircled by the surges of the whole war rather than let the man either fall back or press on --- and flung Mutina, like a curb, upon his exultant fury. And when he had hemmed Mutina round with siege-works and ramparts, and neither the dignity of that most flourishing colony nor the majesty of a consul-designate could deter him from his parricide, then --- and I call you, and the Roman people, and all the gods who watch over this city, to witness --- against my will and over my resistance, three consular envoys were sent to a captain of brigands and gladiators.

\ciceroSection{21} Who was ever so barbarous, so monstrous, so savage? He did not hear them, he did not answer; and not them alone, present before him, but far more us, by whom they had been sent, he scorned and held as nothing. And afterward, what crime, what outrage did the parricide leave undone? He besieges our colonists, an army of the Roman people, a commander, a consul-designate; he ravages the lands of the best of our citizens; the foulest of enemies, he threatens all good men with the cross and the rack. What peace can there be with such a man, Marcus Lepidus? --- a man upon whom the commonwealth seems unable to glut itself with any punishment whatever.

\ciceroSection{22} But if anyone could still doubt that there can be no fellowship between this order and the Roman people and that most savage beast, he will surely cease to doubt once he has heard this letter, which I have just received, sent to me by the consul Hirtius. While I read it out, and while I argue briefly upon its several sentences, I would have you, senators, hear me attentively, as you have done thus far. ``Antonius to Hirtius and Caesar.'' He calls neither himself imperator, nor Hirtius consul, nor Caesar propraetor. Cleverly enough, this: he chose rather to lay aside the title that was not his own than to grant them theirs. ``On learning of the death of Gaius Trebonius I rejoiced no more than I grieved.'' Mark what he says he rejoiced at, and what he grieved at: you will deliberate the more easily about peace. ``That a criminal has paid the penalty to the ashes and bones of an illustrious man, and that the divine power of the gods has shown itself within the turning of the year --- whether the parricide's punishment is now exacted or hangs over him --- is matter for rejoicing.'' O Spartacus! --- for by what name should I rather call you, beside whose unspeakable crimes Catiline seems to have been bearable? You have dared to write that we should rejoice that Trebonius has paid the penalty? Trebonius a criminal? For what crime, save that on the Ides of March he drew you aside from the destruction that was your due?

\ciceroSection{23} Come, you rejoice at this: let us see what you take amiss. ``That Dolabella has been judged by the Senate an enemy of the Roman people because he killed an assassin, and that the son of a buffoon should be seen to be dearer to the commonwealth than Gaius Caesar, the father of his country --- this is matter for groaning.'' Why do you groan? That Dolabella is an enemy? What? Do you not see that you yourself have been judged an enemy --- now that a levy has been raised throughout all Italy, the consuls sent out, Caesar honoured, the soldier's cloak at last assumed? And what is there, you criminal, for you to groan at, that Dolabella has been judged an enemy by the Senate --- you who hold that there is no order at all, and who set yourself this as your reason for waging war: to wipe out the Senate utterly, so that all the other good and wealthy men may follow that highest order down? But he calls Trebonius the son of a buffoon. As though that splendid Roman knight, the father of Trebonius, were unknown to us! And this man dares to despise the low birth of anyone --- he who has had children by Fadia?

\ciceroSection{24} ``But the bitterest thing of all is that you, Aulus Hirtius, adorned with Caesar's favours, and left by him such a man as you yourself marvel to be ---'' I cannot, for my part, deny that Hirtius was adorned by Caesar; but those adornments shine because they rest upon valour and industry. You, however, who cannot deny that you were adorned by the same Caesar --- what would you be, had he not bestowed so much upon you? Whither would your own valour have carried you, whither your birth? You would have spent the whole season of your life in brothels, cookshops, dice, and wine, as you used to do, when you would lay your chin --- and your mind --- in the laps of actresses. ``And you, boy ---'' He calls him `boy,' one whom he has known, and shall know, to be not only a man, but a most valiant man. `Boy' is indeed a name of his age; but it least of all becomes the lips of a man who is offering his own madness to the boy as a path to glory. ``You who owe everything to a name.'' He owes, indeed --- and pays it back nobly.

\ciceroSection{25} For if that man is `the father of his country,' as you call him --- I will look later to what I myself think --- why is this man not the truer father, by whom we surely have our lives, snatched out of your most criminal hands? ``Is this his aim, that Dolabella should be lawfully degraded?'' A base proceeding, indeed --- by which the authority of this most august order is upheld against the madness of a most cruel gladiator! ``And that this poisoner-woman should be freed from the siege?'' Do you dare to call `poisoner' the man who has found the remedies against your poisons? --- a man whom you so besiege, you new Hannibal, or whatever sharper general there ever was, that you besiege your own self, and could not, even if you wished, extricate yourself from where you are. If you draw back, all will pursue you from every side; if you stay, you will be stuck fast. With good reason, surely, you call him `poisoner,' the man by whom you see a present destruction prepared for you. ``So that Cassius and Brutus may be as powerful as possible!''

\ciceroSection{26} You would think it was Censorinus speaking, or Ventidius, or even the Antonii themselves. But why should they be unwilling that men be powerful who are not only the best and noblest, but joined even with us in the defence of the commonwealth? ``No doubt you look on these matters in the same way as on those before.'' On what matters, pray? ``You call Pompey's camp a senate.'' Or should we rather call your camp a senate? --- in which you, forsooth, are the one of consular rank, whose whole consulship has been torn out of every record of the monuments; in which there are two praetors who, without cause, despaired of keeping their standing --- for we are upholding Caesar's grants; among the ex-praetors, Philadelphus Annius and the `blameless' Gallius; among the ex-aediles, Bestia, the punching-bag of my lungs and my voice, and Trebellius, the patron of good faith and the defrauder of his creditors, and Quintus Caelius, a man shattered in his purse, and the very pillar of Antony's friends, Cotyla Varius, whom Antony, for his sport, used to have flogged with thongs at a banquet by the public slaves; the septemvirs Lento and Nucula; then the darling and beloved of the Roman people, Lucius Antonius; first two tribunes-designate: Tullus Hostilius, who by his own right scrawled his name upon the gate by which --- having failed to betray his commander --- he made his exit; the other tribune-designate is one Insteius, a stout brigand, so they say, though they report that, as a bath-keeper at Pisaurum, he was a man of moderation.

\ciceroSection{27} Other ex-tribunes follow, Titus Plancus foremost among them: who, had he loved the Senate, would never have burned down its house. Condemned for that crime, he returned by arms to the city he had left by the law. But this he has in common with many of his sort. One thing, however, is not true in this Plancus, which is wont to be said by way of proverb --- that he cannot perish unless his legs are broken. They have been broken, and he lives. Yet this too, like much else, let it be entered to Aquila's credit. There is also a Decius there, descended, I suppose, from those famous Decii, the `Mice'; and so he has gnawed at Caesar's gifts: the memory of the Decii, indeed, after a long interval, has been revived through this illustrious man. And how can I pass over Saxa Decidius, a fellow fetched from the farthest tribes, that we might see as a tribune of the plebs one we had never seen as a citizen?

\ciceroSection{28} There is, indeed, a second Saserna; but they all have so great a likeness to one another that I go astray among their first names. Nor, in truth, must Extitius, the brother of Philadelphus, the quaestor, be passed over --- lest, if I am silent about so illustrious a young man, I should seem to bear Antony a grudge. There is also a certain Asinius, a volunteer senator, chosen by his own self. He found the Senate-house open after Caesar's death; he changed his shoes; in an instant he was a senator. I do not know Sextus Albesius; but for all that, I have never met anyone so foul-mouthed as to deny him worthy of Antony's senate. I think I have passed some over; yet of those who came to mind I could not keep silent. Relying, then, on such a senate, he despises the Pompeian senate, in which we were ten of consular rank: and had they all been living, this war would never have been at all; daring would have given way before authority.

\ciceroSection{29} But how great a safeguard there was in the rest may be understood from this --- that I, left alone out of many, with your help crushed and broke the daring of the exulting freebooter. And if Fortune had not snatched from us, only lately, Servius Sulpicius, and before him his colleague Marcus Marcellus --- what citizens, what men they were! --- if the commonwealth could have kept the two consuls, dearest friends of the fatherland, driven out of Italy together; if it could have kept Lucius Afranius, that consummate commander; Publius Lentulus, a citizen without peer, both in all else and in the saving of my own life; Marcus Bibulus, whose steadfastness toward the commonwealth has always been deservedly praised; Lucius Domitius, a most excellent citizen; Appius Claudius, endowed with equal nobility and equal goodwill; Publius Scipio, an illustrious man, the very image of his ancestors --- then surely, with such men of consular rank, the Pompeian senate would be no thing to despise.

\ciceroSection{30} Which, then, was the juster, which the better for the commonwealth: that Gnaeus Pompeius should live, or that Antony, the buyer-up of Gnaeus Pompeius's confiscated goods, should live? And the men of praetorian rank! Whose chief was Marcus Cato --- and he the chief of all the nations of the earth in virtue. Why should I recount the other most illustrious men? You know them all. I fear rather to be thought tedious in the listing than ungrateful in the passing-over. What men of aedilician, of tribunician, of quaestorian rank! In short: such was the dignity and the number of those senators that great excuse is needed by those who did not come into that camp. Now attend to the rest. ``You had the defeated Cicero as your leader.'' The more gladly do I hear `leader,' since he certainly says it against his will; for about `defeated' I care nothing. For it is my destiny that, without the commonwealth, I can neither be conquered nor conquer. ``You fortify Macedonia with armies.'' Yes --- and we wrested it from your brother, who in no way falls short of you. ``You committed Africa to Varus, twice captured.'' Here he thinks he is going to law with his brother Gaius. ``You sent Cassius into Syria.'' Do you not feel, then, that to this cause the whole earth lies open, while you have nowhere outside your own siege-works to set the print of your foot? ``You suffered Casca to hold the tribunate.'' What then?

\ciceroSection{31} What then? Were we to drive a Marullus, a Caesetius, out of public life --- through the very man by whom we have secured that neither this same thing nor many of the kind can come to pass hereafter? ``You have stripped the Luperci of the Julian revenues.'' Does he dare make mention of the Luperci? Does he not shudder at the memory of that day on which, sodden with wine, smeared with unguents, naked, he dared to urge the groaning Roman people into slavery? ``You have abolished the veterans' colonies, founded by law and by decree of the Senate.'' Did we abolish them, or did we rather ratify them by a law passed in the centuriate assembly? Take care that you yourself have not undone the veterans --- even those who were already ruined men --- and led them down into a place from which they themselves now feel they will never come out.

\ciceroSection{32} ``You promise to give back to the Massilians what was taken from them by the law of war.'' I make no argument about the law of war --- an argument easier to make than necessary --- but mark this nonetheless, senators: how truly Antony was born the enemy of this commonwealth, who so bitterly hates that city which he knows to have always been the truest friend to this commonwealth. ``You keep saying that no living Pompeian is bound by the Hirtian law.'' Who, pray, now makes any mention of the Hirtian law? Of which I think the proposer himself repents no less than those against whom it was passed. To my mind, indeed, it is not right to call that thing a law at all; and even granting it is a law, we ought not to think it the law of Hirtius. ``You suborned Brutus with the Apuleian money.'' What of it? If the commonwealth had armed so excellent a man with all its resources, what good man, pray, would have repented of it? For without money he could not have maintained an army, nor without an army have captured your brother.

\ciceroSection{33} ``You have praised the beheading of Petraeus and Menedemus, men granted citizenship and guest-friends of Caesar.'' We did not praise it --- we did not even hear of it. For surely, in so great a disturbance of the commonwealth, we had much to think about concerning two worthless little Greeks! ``You have shown no concern that Theopompus, stripped bare, driven out by force by Trebonius, took refuge at Alexandria.'' A grave charge against the Senate! We have neglected Theopompus, that paragon of a man --- who, where on earth he may be, what he is doing, whether in short he is alive or dead, knows or cares? ``You see Servius Galba in the camp, girt with the very same dagger.'' To you I make no answer about Galba, a most brave and steadfast citizen: he will be here in person; he himself, present before you, and that dagger you charge him with, shall make their answer. ``You have drawn together my soldiers, or the veterans, as if for the destruction of those who killed Caesar; and these same men, all unawares, you have driven against the perils of their own quaestor, their commander, their fellow soldiers.'' No doubt we deceived them, played them false: the Martian legion knew nothing, the Fourth, the veterans had no notion of what was afoot; it was not the Senate's authority, not the people's liberty that they were following: they wished to avenge Caesar's death --- which they all believed had been the work of fate --- and wished you, of course, to be safe and blessed and flourishing!

\ciceroSection{34} O wretched man --- both in your case itself, and in this very thing, that you do not feel how wretched you are! But hear the gravest charge of all. ``In a word, what have you not either approved or done that he would do, if he should come back to life'' --- who? for I suppose he will bring forward the example of some scoundrel --- ``Gnaeus Pompeius himself?'' O shame on us, if indeed we have imitated Gnaeus Pompeius! ``Or his son, if only he could?'' He will be able, believe me: for in a few days he will move into both his father's house and his father's gardens. ``Finally, you say that peace cannot be made unless I either release Brutus or aid him with grain.'' Others say so: I, for my part, think that not even if you should do those things would there ever be peace for this city with you. ``What? Does this please those veterans of yours, for whom as yet everything stands untouched?'' I have seen nothing stand so untouched as that they should set about attacking a commander whom they hate with such zeal and such accord.

\ciceroSection{35} ``Whom now you have come to corrupt with flatteries and poisoned gifts.'' Are men corrupted who have been persuaded to hunt down a most foul enemy in a most just war? ``But you are bringing aid to the soldiers shut up within. I have no objection to their being safe and going wherever they please, if only they suffer the man who has deserved it to perish.'' How kind of him! And so, forsooth, the soldiers, having made the most of Antony's generosity, deserted their commander and in their terror went over to the enemy --- whereas, had the matter not rested upon them, Dolabella would have paid the death-offering to his commander no sooner than Antony would have paid his to his colleague as well.

\ciceroSection{36} ``You write that mention has been made of concord in the Senate, and that there are five consular envoys. It is hard to believe that the men who drove me headlong --- when I was offering the fairest terms, and even meaning to relax something of these --- will think to do anything in a moderate or humane spirit. It is scarcely even likely that those who judged Dolabella an enemy for a most upright deed can spare us, who are of the same mind.'' Does he not seem plainly enough to confess that he has entered into a partnership with Dolabella in every crime? Do you not see that from one single spring all his villainies flow? In short, he himself admits --- shrewdly enough, this --- that those who have judged Dolabella an enemy for a most upright deed (so it seems to Antony) cannot spare him, who is of the same mind.

\ciceroSection{37} What are you to do with such a man, who has committed it to writing and to record that he had so agreed with Dolabella that the one should torture and kill Trebonius and, if he could, Brutus and Cassius too, while he himself would inflict the same torments upon us? O what a citizen to be preserved, and by so pious and just a compact! He even complains that his own terms were rejected --- terms fair and modest, forsooth: that he should have Further Gaul, a province most suited for renewing and equipping a war; that the Larks should sit as jurors in the third decury, that is, that there should be a sanctuary for crimes in the foulest dregs of the commonwealth; that his own acts should be ratified --- he of whose consulship not a trace remains. He made provision too for Lucius Antonius, who had been a most even-handed surveyor of land both private and public, with Nucula and Lento for his colleagues.

\ciceroSection{38} ``Therefore consider rather which is the more elegant, and the more useful to the party: to avenge the death of Trebonius, or that of Caesar; and whether it is the fairer course that we should clash, so that the cause of the Pompeians, so often cut down, may the more easily come to life again, or that we should agree together, lest we be a mockery to our enemies.'' If it had been cut down, it would never rise again: which may be the lot of you and yours. ``Which,'' says he, ``is the more elegant.'' And yet in this war it is elegance we are inquiring after!

\ciceroSection{39} ``And the more useful to the party.'' Parties, you madman, are spoken of in the Forum, in the Senate-house. You have taken up an abominable war against the fatherland; you assault Mutina, you beleaguer a consul-designate; against you two consuls wage war, and with them Caesar as propraetor; all Italy is up in arms against you. Do you call that a `party,' rather than a revolt from the Roman people? Are we to avenge the death of Trebonius rather than of Caesar? Trebonius we have avenged enough, with Dolabella judged an enemy; Caesar's death is most easily defended by forgetfulness and silence. But mark what he is contriving. When he holds that Caesar's death must be avenged, he proposes death not only for those who did that deed, but for any too who did not take it ill.

\ciceroSection{40} ``To whom, whichever of us shall fall, it will be a gain --- a spectacle which Fortune herself has so far avoided, that she might not see two battle-lines of one body fighting it out with Cicero as their trainer of gladiators: a man so lucky that he has cheated you with the very same trappings with which he boasts to have cheated Caesar.'' He goes on heaping abuse upon me, as though, forsooth, his earlier abuse had turned out most beautifully --- the man whom I shall hand down, branded with the truest marks of reproach, to the everlasting memory of mankind. I, a trainer of gladiators? And not a foolish one, either: for I want the worse slaughtered and the better to win. ``Whichever shall fall,'' he writes, ``it will be a gain to us.''

\ciceroSection{41} O glorious gain --- when, with you the victor (may the gods avert the omen!), the death of those who shall have departed this life without torture will be a blessed one! He says that Hirtius and Caesar were `cheated by me with the same trappings.' What trapping, pray, have I bestowed on Hirtius up to now? For to Caesar more and greater things are owed. But do you dare to say that the father was cheated by me? You --- you, I say --- killed him at the Lupercalia: whose priesthood, you most ungrateful of men, why have you left undischarged? But now mark the admirable gravity and constancy of this great and illustrious man:

\ciceroSection{42} ``For my part, I am resolved to bear neither insult to myself nor to my own; not to desert the party that Pompey hated, nor to suffer the veterans to be stirred from their settlements, nor to let men be dragged off one by one to torture, nor to break the faith I have given to Dolabella.'' I pass over the rest: the faith pledged to Dolabella, that most holy man --- a pious man cannot desert it. What faith? The faith of slaughtering every best citizen, of partitioning the city and Italy, of laying waste and plundering the provinces? For what else was there to be sealed by treaty and by faith between Antony and Dolabella, those most filthy parricides?

\ciceroSection{43} ``Nor to violate my partnership with Lepidus, a most pious man.'' You, in partnership with Lepidus --- or with anyone, I will not say a good citizen, such as he is, but a man in his right mind? Your aim is to have Lepidus thought either impious or insane. You accomplish nothing --- though about the one it is hard to give assurance --- least of all about Lepidus, whom I shall never fear; I shall hope well of him, as long as I am allowed to. Lepidus wished to call you back from your frenzy, not to be a helper of your madness. And you, moreover, seek out not even the pious, but the `most pious'; and a word that does not exist at all in the Latin tongue, you, for your divine piety's sake, bring in afresh.

\ciceroSection{44} ``Nor to betray Plancus, the sharer of my counsels.'' Plancus your sharer? --- whose memorable and godlike valour brings light to the commonwealth --- unless, perhaps, you imagine he is coming to your aid with his most valiant legions and a great force of Gallic horse and foot --- and who, unless before his arrival you have paid the commonwealth its penalties, will himself bear off the chief honour of this war. For though the first defences are the more useful to the commonwealth, yet the last are the more welcome.

\ciceroSection{45} But now he collects himself, and at the last begins to play the philosopher: ``If the immortal gods, as I hope, shall aid me as I go forward with right understanding, I shall live gladly. But if another fate awaits me, I take in advance my joys of your punishments. For if the Pompeians are so insolent in defeat, what they will be in victory you shall rather find out for yourselves.'' You are welcome to take your joys in advance: for your war is not with the Pompeians, but with the whole commonwealth. All hate you --- gods and men, highest, middling, lowest; citizens and foreigners; men and women; free men and slaves. We felt this lately at a false report; before long we shall feel it at a true one. And if you turn these things over with yourself, you will die with a calmer mind and a greater consolation.

\ciceroSection{46} ``In a word, the sum of my judgment comes to this: that I can endure the wrongs done to my own people, if the doers are willing either to forget that they did them, or stand ready, together with us, to avenge the death of Caesar.'' Now that this sentiment of Antony's is known, do you suppose that the consuls Aulus Hirtius or Gaius Pansa will hesitate to go over to Antony, to besiege Brutus, to long to storm Mutina? Why do I speak of Pansa and Hirtius? Caesar, a young man of singular devotion --- will he be able to hold himself back from pursuing his father's vengeance in the blood of Decimus Brutus? And so they brought it about that, on the reading of this letter, they drew nearer to the siege-works. The greater, then, is the young Caesar, and born by the greater favour of the immortal gods for the commonwealth, in that he has never been led astray by any show of his father's name nor by any filial piety, and understands that the greatest piety is contained in the preservation of the fatherland.

\ciceroSection{47} But if it were a contest of parties --- whose very name is utterly extinct --- would Antony and Ventidius be the fitter defenders of Caesar's party than, first, Caesar, a young man of the highest devotion and mindfulness of his parent, and then Pansa and Hirtius, who held, as it were, the two wings of Caesar at the time when those were truly called parties? But what manner of parties are these, when on the one side stand the authority of the Senate, the liberty of the Roman people, the safety of the commonwealth, and on the other the slaughter of good men, the partitioning of the city and of Italy? Let us come at last to the close. ``I do not believe the envoys will come.'' He knows me well: let the rest come --- especially now that the example of Dolabella has been set before them. Will the envoys, I suppose, be protected by a more sacred law than the two consuls against whom he bears arms, than Caesar whose father's priest he is, than the consul-designate whom he assaults, than Mutina which he besieges, than the fatherland which he threatens with fire and sword?

\ciceroSection{48} ``When they have come, I shall learn what they demand.'' Why do you not betake yourself to a foul plague and a foul torment? Is anyone to come to you but one like Ventidius? We sent the foremost of men to quench the rising blaze; you rejected them: are we now to send them into so great a fire, and one so deeply set, when you have left yourself no room not only for peace, but not even for surrender? This letter, senators, I have read out, not because I thought him worthy of it, but that by his own confessions you might see all his parricides laid open to view.

\ciceroSection{49} Would Marcus Lepidus --- a man most richly endowed with all the goods both of valour and of fortune --- wish for peace with this man, or think it possible, if he saw all this? ``Sooner shall fire mingle with the waves,'' as some poet or other says --- sooner, in short, shall all things else come to pass, than that either the commonwealth and the Antonii, or the Antonii and the commonwealth, should be reconciled. These men are monsters, portents, prodigies of the commonwealth. It were better for this city to be moved from its foundations, and to migrate, if it could, into other lands, where it should hear `neither the deeds nor the name' of the Antonii, than to see them within these walls --- men cast out by Caesar's valour, held off by Brutus's. To conquer is most to be wished for; the next best is to think no fate too hard to bear for the dignity and the liberty of the fatherland. What remains is not the third, but the last of all: to take upon oneself the greatest disgrace out of love of life.

\ciceroSection{50} Since this is so, concerning the instructions and the dispatch of Marcus Lepidus, that most illustrious man, I assent to Servilius; and I move this further: that Magnus Pompeius, son of Gnaeus, in accordance with the spirit and the zeal of his father and his ancestors toward the commonwealth, and with his own old valour, industry, and goodwill, has done well in promising to the Senate and the Roman people his own service and that of those whom he had with him; that this is welcome and acceptable to the Senate and the Roman people; and that it shall be to his honour and his dignity. This may either be joined with the present decree of the Senate, or it can be kept separate and written out apart, so that Pompeius may be seen to have been commended by a decree of the Senate of his own.

